This section motivates reinforcement learning (RL) for large language models (LLMs), defines scope, and previews the survey roadmap. We open with an executive overview, then explain why RL became central to LLM post-training, then state our scope and notation, and finally compare this survey with prior work. Representative production systems shaping our scope include: InstructGPT (2022, RLHF on 175B GPT-3 base), ChatGPT (2022, first mass-deployed RLHF chat), Claude (2022, Constitutional-AI RLAIF), Sparrow (2022, rule-augmented RLHF), Llama 2-Chat (2023, dual helpfulness/safety RMs), Llama 3-Instruct (2024, iterative DPO at 405B), Qwen2.5-Math (2024, RL on synthetic chain-of-thought), DeepSeek-V3 (2024, 671B MoE base), DeepSeek-R1 (2025, GRPO+RLVR \emph{Nature} paper), OpenAI o1 (2024, RL on reasoning traces), and OpenAI o3 (2025, deeper test-time RL). The survey returns to each of these systems in the relevant later sections. \textbf{Executive overview.} Reinforcement learning (RL) is the dominant post-training stage for large language models. It converts pretrained next-token predictors into deployable assistants, reasoners, coders, and agents. This survey covers four years of consolidation across four threads: (i) chat alignment via RLHF (InstructGPT, Claude, Llama 2/3); (ii) verifiable-reward reasoning via RLVR (DeepSeek-R1, OpenAI o1, Qwen2.5-Math); (iii) offline preference optimization via the DPO family (DPO, IPO, KTO, ORPO, SimPO); and (iv) agentic and multimodal RL (WebGPT, \ldots{}
